% !PRINTED 569
\weberpage{569}

%<<CONT
worin $\nu$ ebenso wie $\mu$ eine ganze Zahl des Körpers $\Omega$ ist, die \gesperrt{nicht durch $\mu$ teilbar ist.}

Ist diese Bedingung erfüllt, so ist $g$ eine Wurzel von $P_\mu(x)$. Weil aber $\wp(\mu u)$ in der $u$-Ebene für $u\equiv 0\ (\mathrm{modd}.\ \omega_1,\omega_2)$ unendlich in der zweiten Ordnung wird, so ist $P_\mu(x)$ durch $(x-g)^2$ teilbar, es sei denn, daß
\[ \wp(u) - g \]
für $u=\nu:\mu$ selbst in der zweiten Ordnung verschwindet; dies tritt wegen § 46, (18) dann ein, wenn
\[ \wp'\left(\dfrac{\nu}{\mu}\right) = 0 \]
ist, also wenn $\nu:\mu$ eine halbe Periode ist, und also
\weq{6}{g = e_1, e_2, e_3}
wird.

\begin{quote}
{\small 1. \gesperrt{Nach (4) ist, wenn $\omega_1=1$, $\omega_2=\omega$ gesetzt ist, jede ganze Zahl des Körpers $\Omega$ eine Periode von $\wp(u)$. Dies gilt nicht umgekehrt; da aber $A\omega$ eine ganze Zahl in $\Omega$ ist, so muß jede Periode durch Multiplikation mit $A$ in eine ganze Zahl verwandelt werden.}}
\end{quote}

Daraus folgt, daß zwei Werte
\[ g = \wp\left(\dfrac{\nu}{\mu}\right), \qquad g' = \wp\left(\dfrac{\nu'}{\mu}\right) \]
dann und nur dann einander gleich sind, wenn
\[ \nu \equiv \pm \nu' \pmod{\mu} \]
ist. Denn ist $g=g'$, so muß entweder
\[ \dfrac{\nu+\nu'}{\mu} \quad \text{oder} \quad \dfrac{\nu-\nu'}{\mu} \]
durch Multiplikation mit $A$ in eine ganze Zahl verwandelt werden. $A$ ist aber relativ prim zu $\mu$, und folglich muß eine dieser beiden Zahlen selbst ganz sein. Die Anzahl der inkongruenten Werte von $\nu$ ist aber gleich $N(\mu) = m$ (Bd.\ II, § 165).

Lassen wir also $\nu$ ein volles Restsystem nach dem Modul $\mu$ mit Ausschluß der Null durchlaufen, so bekommen wir jeden Wert von $g$ zweimal, außer wenn
\weq{7}{2\nu \equiv 0 \pmod{\mu},}
und in diesem Falle ist $\wp\left(\dfrac{\nu}{\mu}\right)$ einer der Werte $e$. Daraus folgt:
%CONT>>
